\documentclass{beamer}
\usepackage{ctex}
\usepackage{tikz}
\usepackage{tikz-cd}
\usepackage{xypic}

\usepackage{unicode-math}
%\setmathfont{STIX Two Math}
%\setmathfont{TeX Gyre Termes Math}
%\setmathfont{Libertinus Math}
%\setmathfont{Fira Math}

% MACROS

\newcommand{\NN}{\mathbb{N}} % natural numbers
\newcommand{\RR}{\mathbb{R}} % real numbers
\newcommand{\ZZ}{\mathbb{Z}} % integers

\newcommand{\theory}[1]{\mathsf{#1}} % A named theory
\newcommand{\signature}[1]{\Sigma_{\theory{#1}}} % The signature of a theory
\newcommand{\equations}[1]{\mathcal{E}_{\theory{#1}}} % Equations of a theory

\newcommand{\Mod}[1]{\text{\textbf{Mod}}(\theory{#1})} % Models of a theory in Set
\newcommand{\ModC}[2]{\text{\textbf{Mod}}_{#1}(\theory{#2})} % Models of a theory in a category
\newcommand{\ComodC}[2]{\text{\textbf{Comod}}_{#1}(\theory{#2})} % Comodels of a theory in a category

\newcommand{\category}[1]{\text{\textbf{#1}}} % A named category
\newcommand{\Set}{\category{Set}} % The category of sets
\newcommand{\opcat}[1]{{#1}^{\mathsf{op}}} % Opposite category

\newcommand{\carrier}[1]{|#1|} % carrier of a model
\newcommand{\Free}[2]{F_{\theory{#1}}(#2)} % free model
\newcommand{\FreeFun}[1]{F_{\theory{#1}}} % free model functor

\newcommand{\all}[1]{\forall #1 \,.\,} % universal quantifier
\newcommand{\some}[1]{\exists #1 \,.\,} % existential quantifier
\newcommand{\set}[1]{\{#1\}} % set description
\newcommand{\such}{∣}

\newcommand{\lam}[1]{\lambda #1 \,.\,}

\newcommand{\family}[2]{\{#1\}_{#2}} % a family

\newcommand{\finpow}[1]{\mathcal{P}_{{<}\omega}(#1)} % finite powerset

\newcommand{\Tree}[2]{\mathsf{Tree}_{#1}(#2)} % trees over a signature
\newcommand{\leaf}[1]{\return{#1}} % the embedding of generators into trees

\newcommand{\op}[1]{\mathsf{op}_{#1}} % an operation symbol
\newcommand{\arity}[1]{\mathsf{ar}_{#1}} % arity of a symbol
\newcommand{\opdecl}[3]{#1 : #2 \leadsto #3} % operation declaration

\newcommand{\one}{\mathsf{1}} % the terminal object
\newcommand{\unit}{{\text{\small$()$}}} % the sole element of the singleton

\newcommand{\lift}[1]{#1^\dagger} % monad lift

\newcommand{\Cinfty}{\mathcal{C}^\infty}

\newcommand{\sem}[1]{[\![#1]\!]} % semantic bracket

\newcommand{\bool}{\mathsf{bool}} % booleans
\newcommand{\true}{\mathsf{true}}
\newcommand{\false}{\mathsf{false}}
\newcommand{\cond}[3]{\mathsf{if}\;#1\;\mathsf{then}\;#2\;\mathsf{else}\;#3}

\newcommand{\hto}{\Rightarrow} % Handler arrow

\newcommand{\defeq}{\mathbin{{:}{=}}} % definitional equality

\newcommand{\tensor}[2]{#1 \otimes #2} % tensor of a model and a comodel
\newcommand{\opair}[2]{\langle #1, #2 \rangle} % the tensor of a progrqm and a world

\newcommand{\ct}[1]{\underline{#1}} % computation type symbol
\newcommand{\dirt}[2]{#1 \mathbin{!} #2} % type with dirt

%%% Macros for the programming language

% General syntactic constructs
\newcommand{\kode}[1]{\mathsf{#1}}

% Core syntax
\newcommand{\seq}[2]{\kode{do}\; #1 \leftarrow #2 \;\kode{in}\;}
\newcommand{\handler}{\kode{handler}\;}
\newcommand{\xopgen}[2]{\overline{#1}(#2)}
\newcommand{\opgen}[2]{\xopgen{\kode{#1}}{#2}}
\newcommand{\opcall}[3]{\kode{#1}(#2, #3)}
\newcommand{\opclause}[3]{#1(#2; #3) \mapsto}
\newcommand{\return}[1]{\kode{return}\;#1}
\newcommand{\retclause}[1]{\return{#1} \mapsto}
\newcommand{\withhandle}[2]{\kode{with}\; #1\; \kode{handle}\; #2}

\title{Algebraic effects}

\date{}

\begin{document}

\begin{frame}
    % Print the title page as the first slide
    \titlepage
\end{frame}

\begin{frame}{签名 $\Sigma$}
  $\family{(\op{i}, \arity{i})}{i}$

  \vspace{6mm}

  幺半群(monoid)对应的签名为 $\{(\mathsf{u}, 0), (\mathsf{m}, 2)\}$
\end{frame}

\begin{frame}{$\Sigma$-项(term)}
 不同变量$x_1, …, x_k$的一个列表(可能为空)被称为上下文(context)

 \vspace{6mm}

 在上下文$x_1, …, x_k$中的$\Sigma$-项根据以下规则归纳定义
 \begin{enumerate}
  \item 在上下文$x_1, …, x_k$中, 每一个变量 $x_i$ 都是一个 $\Sigma$-项
  \item 如果$t_1, …, t_{\arity{i}}$是上下文$x_1, …, x_k$中的$\Sigma$-项， 那么 $\op{i}(t_1, …, t_{\arity{i}})$ 是上下文 $x_1, …, x_k$中的一个$\Sigma$-项
 \end{enumerate}

 \vspace{6mm}

 我们用
 \begin{equation*}
  x_1, …, x_k ∣ t
\end{equation*}
表示$t$是上下文$x_1, …, x_k$中的一个$\Sigma$-项

\vspace{6mm}

空上下文中的$\Sigma$-项被称为闭(closed)$\Sigma$-项. 显然, 闭$\Sigma$-项中不含变量
\end{frame}

\begin{frame}{$\Sigma$-等式}
  一个$\Sigma$-等式就是上下文$x_1, …, x_k$中的一对$\Sigma$-项$\ell$和$r$. 一般记为
\begin{equation*}
  x_1, …, x_k ∣ ℓ = r
\end{equation*}
\end{frame}

\begin{frame}{代数理论(Algebraic theories)}
  一个代数理论 $\theory{T} = (\signature{T}, \equations{T})$, 也被叫作一个等式理论(equational theory), 其中$\signature{T}$是一个签名, $\equations{T}$ 是一组 $\signature{T}$-等式.
\end{frame}

\begin{frame}{签名的解释}
给定一个签名 $\Sigma$. $\Sigma$ 的一个解释 $I$ 由以下数据构成:
\begin{enumerate}
  \item 一个集合 $\carrier{I}$, 叫作承载集(carrier).
  \item 对每一个操作符 $\op{i}$, 有一个映射
    \begin{equation*}
    \sem{\op{i}}_I : \underbrace{\carrier{I} × ⋯ × \carrier{I}}_{\arity{i}} → \carrier{I},
  \end{equation*}
  称为运算.
\end{enumerate}
\end{frame}

\begin{frame}{$\Sigma$-项的解释}
  签名 $\Sigma$ 的一个解释 $I$ 可以被扩展到 $\Sigma$-项. 上下文中的一个 $\Sigma$-项
\begin{equation*}
  x_1, …, x_k ∣ t
\end{equation*}
由一个映射
\begin{equation*}
  \sem{x_1, …, x_k ∣ t}_I : \carrier{I}^k → \carrier{I},
\end{equation*}
来解释. 定义如下:
\begin{enumerate}
  \item 变量$x_i$被解释为第$i$个投影
    \begin{equation*}
    \sem{x_1, …, x_k ∣  x_i}_I = \pi_i : \carrier{I}^k → \carrier{I},
  \end{equation*}
  \item 上下文中的一个复合项
  \begin{equation*}
    x_1, …, x_k ∣ \op{i}(t_1, …, t_{\arity{i}})
  \end{equation*}
  被解释为映射的复合
    \begin{equation*}
    \xymatrix@+6em{
      {\carrier{I}^k} \ar[r]^{(\sem{t_1}_I, …, \sem{t_{\arity{i}}}_I)}
      &
      {\carrier{I}^{\arity{i}}} \ar[r]^{\sem{\op{i}}_I}
      &
      {\carrier{I}}
    }
  \end{equation*}
  (为了简单起见, 我们在合适的地方省略了上下文$x_1, …, x_k$)
\end{enumerate}
\end{frame}

\begin{frame}{代数理论的模型}
  代数理论$\theory{T}$的一个模型$M$, 是签名$\signature{T}$的一个解释, 并且这个解释使得$\equations{T}$中的所有等式成立. 也就是说, 对$\equations{T}$中的任意等式
  \begin{equation*}
  x_1, …, x_k ∣ ℓ = r
\end{equation*}
映射
\begin{equation*}
  \sem{x_1, …, x_k ∣ \ell}_M : \carrier{M}^k → \carrier{M}
  \quad\text{和}\quad
  \sem{x_1, …, x_k ∣ r}_M : \carrier{M}^k → \carrier{M}
\end{equation*}
相等. 我们把$\theory{T}$的模型称为$\theory{T}$-模型或$\theory{T}$-代数.

\end{frame}

\begin{frame}{例}
  每一个代数理论都有一个平凡模型, 它的承载集是单点集$\one$, 它的操作被解释为唯一的映射$\one^k → \one$. 所有等式成立, 因为任意两个映射$\one^k → \one$是相等的.

  \vspace{6mm}

  这就解释了为什么对一个环(ring)来说, 我们不应该要求 $0 ≠ 1$. 因为如果这样, 环的理论就不是一个代数理论.
\end{frame}

\begin{frame}{模型的积}
  假设$L$和$M$是理论$\theory{T}$的两个模型. 我们可以定义模型的积$L× M$如下: 承载集为模型承载集的笛卡尔积
  \begin{equation*}
  \carrier{L × M} = \carrier{L} × \carrier{M},
  \end{equation*}
操作逐点计算
\begin{equation*}
  \sem{\op{i}}_{M × L}(a, b) = (\sem{\op{i}}_M(a), \sem{\op{i}}_L(b)).
\end{equation*}
等式$\equations{T}$在$L× M$上成立, 因为它们在每个坐标分量处分别成立. 此处的积操作可以扩展为任意个数模型的积, 以至于无穷个模型的积.
\end{frame}

\begin{frame}{模型之间的同态}
  令 $L$ 和 $M$ 是代数理论 $\theory{T}$ 的模型. 从 $L$ 到 $M$ 的一个 $\theory{T}$-同态是承载集之间的一个映射 $\phi : \carrier{L} → \carrier{M}$, 并且它和 $\theory{T}$ 的每一个操作符 $\op{i}$ 都是交换的
  \begin{equation*}
  ϕ ∘ \sem{\op{i}}_L = \sem{\op{i}}_M ∘ \underbrace{(\phi, …, \phi)}_{\arity{i}}.
  \end{equation*}

  \vspace{10mm}

  我们可以把代数理论 $\theory{T}$ 的模型组织为一个范畴 $\Mod{T}$: 它的对象是 $\theory{T}$ 的模型, 态射是模型之间的同态.
\end{frame}

\begin{frame}{范畴中的模型}
前面我们考虑的是代数理论在集合上的模型. 更一般地, 我们可以考虑具有有限积的范畴 $\category{C}$ 中的模型. 前面关于解释和模型的定义可以直接转写到 $\category{C}$ 上. $\category{C}$ 中的一个解释 $I$ 定义为
\begin{enumerate}
  \item $\category{C}$ 中的一个对象 $\carrier{I}$, 叫作承载对象(carrier)
  \item 对每一个操作符 $\op{i}$, 有一个$\category{C}$中的态射
    \begin{equation*}
    \sem{\op{i}}_I : \underbrace{\carrier{I} × ⋯ × \carrier{I}}_{\arity{i}} → \carrier{I}.
  \end{equation*}
\end{enumerate}

解释 $I$ 可以扩展到上下文中的 $\Sigma$-项:
\begin{enumerate}
  \item 变量 $x_1, …, x_k ∣ x_i$ 被解释为第 $i$-个投影
    \begin{equation*}
    \sem{x_1, …, x_k ∣  x_i}_I = \pi_i : \carrier{I}^k → \carrier{I},
    \end{equation*}
  \item 上下文中的一个复合项
  \begin{equation*}
    x_1, …, x_k ∣ \op{i}(t_1, …, t_{\arity{i}})
  \end{equation*}
  被解释为态射的复合
  \begin{equation*}
    \xymatrix@+6em{
      {\carrier{I}^k} \ar[r]^{(\sem{t_1}_I, …, \sem{t_{\arity{i}}}_I)}
      &
      {\carrier{I}^{\arity{i}}} \ar[r]^{\sem{\op{i}}_I}
      &
      {\carrier{I}}
    }
  \end{equation*}
\end{enumerate}
\end{frame}

\begin{frame}
  代数理论$\theory{T}$在范畴$\category{C}$中的一个模型, 是签名$\signature{T}$的一个解释$M$, 并且这个解释使得$\equations{T}$中的所有等式成立. 也就是说, 对$\equations{T}$中的任意等式
  \begin{equation*}
  x_1, …, x_k ∣ ℓ = r
\end{equation*}
态射
\begin{equation*}
  \sem{x_1, …, x_k ∣ \ell}_M : \carrier{M}^k → \carrier{M}
  \quad\text{和}\quad
  \sem{x_1, …, x_k ∣ r}_M : \carrier{M}^k → \carrier{M}
\end{equation*}
相等.
\end{frame}

\begin{frame}
同态的定义同样可以扩展到范畴中的模型. 范畴$\category{C}$中的$\theory{T}$-模型 $L$和$M$之间的一个$\theory{T}$-同态是 $\category{C}$ 中的一个态射 $\phi : \carrier{L} \to \carrier{M}$, 使得对$\theory{T}$中的每一个操作符$\op{i}$, $\phi$ 和 $\op{i}$ 的解释交换
\begin{equation*}
  ϕ ∘ \sem{\op{i}}_L = \sem{\op{i}}_M ∘  \underbrace{(\phi, …, \phi)}_{\arity{i}}.
\end{equation*}

范畴$\category{C}$中的$\theory{T}$-模型和$\theory{T}$-同态构成一个范畴$\ModC{\category{C}}{T}$.
\end{frame}

\end{document}